% --- IMPORTS ---
\documentclass[leqno]{article}
\usepackage{graphicx}
\usepackage{dsfont}
\usepackage{mathtools}
\usepackage{amssymb}
\usepackage{amsmath}
\usepackage{amsthm}
\usepackage{float}
\usepackage{ragged2e}
\usepackage{tensor}
\usepackage{fancyhdr}
\usepackage{empheq}
\usepackage[most]{tcolorbox}
\usepackage{enumitem}
\usepackage{dirtytalk}
\usepackage{marginnote}
\usepackage{microtype}
\usepackage[
  a4paper,
  left=0.8in,
  right=1in,
  top=1in,
  bottom=0.5in,
  headheight=14pt,
  headsep=0.4in
]{geometry}
\setlength{\parindent}{0cm}
% --- TikZ Packages ---
\usepackage{pgfplots}
\pgfplotsset{compat=1.18}
\usepgfplotslibrary{fillbetween, polar}
\usetikzlibrary{calc, positioning}
\usetikzlibrary{angles, quotes}
\usetikzlibrary{arrows.meta}
\usetikzlibrary{decorations.pathreplacing, decorations.markings}
\usetikzlibrary{patterns}
\usetikzlibrary{intersections}
% --- URL Packages ---
\usepackage[colorlinks=true,linkcolor=red,citecolor=magenta,urlcolor=black]{hyperref}%
\urlstyle{same}
% --- END OF IMPORTS ---

% --- CUSTOM COMMANDS ---
\RenewDocumentCommand{\marks}{o m}{%
  \IfValueTF{#1}%
    {\marginnote{\raggedright\textbf{[#2]}}[#1]}%
    {\marginnote{\raggedright\textbf{[#2]}}}%
}
\newcommand{\vspacerule}[2]{%
  \par\vspace*{#1}%
  \noindent\hrulefill\par
  \vspace*{#2}%
}
\definecolor{scarlet}{RGB}{205, 40, 75}
\newcommand{\questionitem}{%
  \item\phantomsection
  \addcontentsline{toc}{section}{Question \number\value{enumi}}%
}
\renewcommand{\contentsname}{Questions}
% --- END OF CUSTOM COMMANDS ---

% --- HEADER CONFIG ---
\pagestyle{fancy}
\fancyhead{}
\fancyfoot{}
\renewcommand{\headrulewidth}{0.9pt}
\lhead{\textit{Solutions} - Geometry of Plane Intersections}
\chead{}
\rhead{\href{https://danieljsmith.org}{danieljsmith.org}}
% --- END OF HEADER CONFIG ---

\begin{document}
\vspace*{20pt}
\tableofcontents
\newpage
\begin{enumerate}[
  label=\textbf{\arabic*.},
  leftmargin=*,
  topsep=0pt,
  partopsep=0pt,
  parsep=0pt
]

% ---- Question 1 ----
\questionitem

Three planes have equations
\[
\begin{aligned}
x-y+2z&=3\\
2x+y+z&=4\\
3x+(m+1)z&=7
\end{aligned}
\] \\
where $m$ is a constant. \\
\begin{enumerate}[label=(\roman*)]
  \item Investigate the arrangement of the planes:

  when $m=2$;

  when $m\ne 2$. \marks{6} \\
  \item A student claims that the position vectors $\mathbf{i}-\mathbf{j}+2\mathbf{k}$, $2\mathbf{i}+\mathbf{j}+\mathbf{k}$ and $3\mathbf{i}+4\mathbf{k}$ all lie in a plane through the origin. Determine whether or not the student is correct. \marks{2} 
\end{enumerate}

\vspace*{10pt}

\begin{tcolorbox}[colback=white, colframe=scarlet, title={\textbf{Solution}}, breakable, grow to left by=6mm, grow to right by=10mm]
\begin{enumerate}[label=\textbf{(\roman*)}]
  \item Let the planes be \(P_1,P_2,P_3\), with coefficient matrix
\[
A=\begin{pmatrix}
1 & -1 & 2\\
2 & 1 & 1\\
3 & 0 & m+1
\end{pmatrix}
\]

Its determinant is
\begin{align*}
\det A
&=1\begin{vmatrix}1&1\\0&m+1\end{vmatrix}
-(-1)\begin{vmatrix}2&1\\3&m+1\end{vmatrix}
+2\begin{vmatrix}2&1\\3&0\end{vmatrix}\\
&=1\bigl(1(m+1)-0\bigr)
+1\bigl(2(m+1)-3\bigr)
+2\bigl(0-3\bigr)\\
&=3(m-2)
\end{align*}

If \(m\ne 2\), then \(\det A\ne 0\), so the three planes meet in exactly one point.

To find this point, solve the first two equations:
\begin{align*}
x-y+2z&=3\\
2x+y+z&=4
\end{align*}

Adding gives
\begin{align*}
3x+3z&=7\\
x+z&=\frac73\\
x&=\frac73-z
\end{align*}

Substitute into \(2x+y+z=4\):
\begin{align*}
2\left(\frac73-z\right)+y+z&=4\\
\frac{14}{3}-2z+y+z&=4\\
y-z&=4-\frac{14}{3}\\
y-z&=-\frac23\\
y&=z-\frac23
\end{align*}

Now substitute into the third plane:
\begin{align*}
3x+(m+1)z&=7\\
3\left(\frac73-z\right)+(m+1)z&=7\\
7-3z+(m+1)z&=7\\
(m-2)z&=0
\end{align*}

Since \(m\ne 2\), it follows that \(z=0\). Then
\[
x=\frac73,\qquad y=-\frac23
\]

So when \(m\ne 2\), the three planes intersect at the single point
\[
\left(\frac73,-\frac23,0\right)
\]

If \(m=2\), then the third plane becomes
\[
3x+3z=7
\]

But adding the first two plane equations gives
\begin{align*}
(x-y+2z)+(2x+y+z)&=3+4\\
3x+3z&=7
\end{align*}

So \(P_3\) is obtained from \(P_1\) and \(P_2\), which means the line of intersection of \(P_1\) and \(P_2\) also lies in \(P_3\).

To find that line, let
\[
z=t
\]

Then from \(x+z=\frac73\),
\[
x=\frac73-t
\]

And from \(y=z-\frac23\),
\[
y=t-\frac23
\]

Hence the common line is
\[
x=\frac73-t,\qquad y=t-\frac23,\qquad z=t
\]

Therefore, when \(m=2\), all three planes share the line
\[
\left(\frac73-t,\ t-\frac23,\ t\right)
\]

So the arrangements are:
\begin{itemize}
\item if \(m\ne 2\), the three planes meet at one point \(\left(\frac73,-\frac23,0\right)\)
\item if \(m=2\), the three planes have a common line \(x=\frac73-t,\ y=t-\frac23,\ z=t\), so form a \textbf{sheaf}.
\end{itemize}

  \item The three position vectors are
\[
\begin{pmatrix}1\\-1\\2\end{pmatrix},
\qquad
\begin{pmatrix}2\\1\\1\end{pmatrix},
\qquad
\begin{pmatrix}3\\0\\4\end{pmatrix}
\]

They lie in a plane through the origin if and only if they are coplanar, which is equivalent to their scalar triple product being \(0\). So we test
\begin{align*}
\begin{vmatrix}
1 & -1 & 2\\
2 & 1 & 1\\
3 & 0 & 4
\end{vmatrix}
&=1\begin{vmatrix}1&1\\0&4\end{vmatrix}
-(-1)\begin{vmatrix}2&1\\3&4\end{vmatrix}
+2\begin{vmatrix}2&1\\3&0\end{vmatrix}\\
&=1(4)+1(8-3)+2(0-3) = 3
\end{align*}

Since this determinant is not \(0\), the vectors are not coplanar.

Therefore the student is incorrect.
\end{enumerate}
\end{tcolorbox}


\newpage

% ---- Question 2 ----
\questionitem

\begin{enumerate}[label=\textbf{(\alph*)}]
  \item Consider the linear system $A\mathbf{x}=\mathbf{b}$, where
  \[
  A=\begin{pmatrix}
  3&-2&1\\
  2&1&-4\\
  7&0&-7
  \end{pmatrix},
  \qquad
  \mathbf{x}=\begin{pmatrix}x\\y\\z\end{pmatrix},
  \qquad
  \mathbf{b}=\begin{pmatrix}5\\1\\7\end{pmatrix}
  \]
  Show that $A$ is singular, and deduce that the system does not have a unique solution. \marks{2}\\

  \item Show that the system is consistent. Interpret the solution set geometrically. \marks{3} \\
\end{enumerate}

\vspace*{30pt}

\begin{tcolorbox}[colback=white, colframe=scarlet, title={\textbf{Solution}}, breakable, grow to left by=6mm, grow to right by=10mm]
\begin{enumerate}[label=\textbf{(\alph*)}]
  \item Using a calculator:
\[
\det(A)=0
\]
Therefore \(A\) is singular.

Since a unique solution to \(A\mathbf{x}=\mathbf{b}\) would require \(A\) to be non-singular, the system does not have a unique solution.

\item The system is
\begin{align*}
3x-2y+z&=5 \\
2x+y-4z&=1 \\
7x-7z&=7
\end{align*}

First check consistency. Take the first equation and add twice the second:
\begin{align*}
(3x-2y+z)+2(2x+y-4z)&=5+2(1) \\
7x-7z&=7
\end{align*}
This is exactly the third equation, so the system is consistent.

Now solve using the first two equations. From
\[
2x+y-4z=1
\]
we get
\[
y=1-2x+4z
\]

Substitute into \(3x-2y+z=5\):
\begin{align*}
3x-2(1-2x+4z)+z&=5 \\
3x-2+4x-8z+z&=5 \\
7x-7z-2&=5 \\
7x-7z&=7 \\
x-z&=1 \\
x&=z+1
\end{align*}

Then
\begin{align*}
y&=1-2x+4z \\
&=1-2(z+1)+4z \\
&=1-2z-2+4z \\
&=-1+2z
\end{align*}

Let \(z=t\), where \(t\in\mathbb{R}\). Then
\[
x=1+t,\qquad y=-1+2t,\qquad z=t
\]
So the general solution is
\[
(x,y,z)=(1+t,-1+2t,t),\qquad t\in\mathbb{R}
\]
or equivalently
\[
\begin{pmatrix}x\\y\\z\end{pmatrix}
=
\begin{pmatrix}1\\-1\\0\end{pmatrix}
+t\begin{pmatrix}1\\2\\1\end{pmatrix},
\qquad t\in\mathbb{R}
\]

Geometrically, each equation represents a plane in \(\mathbb{R}^3\). Since the third equation is dependent on the first two and is compatible with them, all three planes intersect in the line
\[
\mathbf{r}=
\begin{pmatrix}1\\-1\\0\end{pmatrix}
+t\begin{pmatrix}1\\2\\1\end{pmatrix},
\qquad t\in\mathbb{R}
\]

Hence the solution set is a line in \(\mathbb{R}^3\), so the planes form a \textbf{sheaf}.
\end{enumerate}
\end{tcolorbox}


\newpage

% ---- Question 3 ----
\questionitem

\begin{enumerate}[label=\textbf{(\alph*)}]
  \item Show that the system of equations
  \[
  \begin{aligned}
  x+2y-z&=3\\
  2x-y+4z&=1\\
  3x+y+3z&=a
  \end{aligned}
  \]
  where $a$ is a constant, does not have a unique solution. \marks{2}\\
  \item Given that $a=4$, show that the system of equations in part (a) is consistent. Interpret this situation geometrically. \marks{3} \\
  \item Given instead that $a\ne 4$, show that the system of equations in part (a) is inconsistent. Interpret this situation geometrically. \marks{2} \\
\end{enumerate}

\vspace*{30pt}

\begin{tcolorbox}[colback=white, colframe=scarlet, title={\textbf{Solution}}, breakable, grow to left by=6mm, grow to right by=10mm]
\begin{enumerate}[label=\textbf{(\alph*)}]
  \item Let the coefficient matrix be
\[
A=\begin{pmatrix}
1&2&-1\\
2&-1&4\\
3&1&3
\end{pmatrix}
\]
and calculate using, for example, a calculator:
\[
\det(A)=0
\]
Therefore $A$ is not invertible, so the system cannot have a unique solution for any value of $a$.

  \item When $a=4$, the third equation is
\[
3x+y+3z=4
\]
Now add the first two equations:
\[
(x+2y-z)+(2x-y+4z)=3+1
\]
which gives
\[
3x+y+3z=4
\]
So the third equation is exactly the sum of the first two, and adds no new condition. Hence the system is consistent.

To find the solutions, solve the first two equations only. Let
\[
z=5t
\]
where $t\in\mathbb{R}$.

From
\[
x+2y-z=3
\]
we have
\[
x=3-2y+z
\]
Substitute into the second equation:
\begin{align*}
2x-y+4z&=1\\
2(3-2y+z)-y+4z&=1\\
6-4y+2z-y+4z&=1\\
6-5y+6z&=1\\
-5y+6z&=-5\\
5y-6z&=5\\
y&=1+\frac{6}{5}z
\end{align*}
Since $z=5t$, this gives
\[
y=1+6t
\]
Then
\begin{align*}
x&=3-2y+z\\
&=3-2(1+6t)+5t\\
&=3-2-12t+5t\\
&=1-7t
\end{align*}
Hence the solutions are
\[
x=1-7t,\qquad y=1+6t,\qquad z=5t,\qquad t\in\mathbb{R}
\]
So there are infinitely many solutions.

Geometrically, each equation represents a plane in $\mathbb{R}^3$. The first two planes intersect in a line, and the third plane contains that same line, so the planes form a \textbf{sheaf}.

  \item If $a\ne 4$, then any point satisfying the first two equations must also satisfy
\begin{align*}
3x+y+3z&=(x+2y-z)+(2x-y+4z)\\
&=3+1\\
&=4
\end{align*}
So every common solution of the first two equations lies on the plane
\[
3x+y+3z=4
\]
It cannot also satisfy
\[
3x+y+3z=a
\]
when $a\ne 4$.

Therefore the system is inconsistent, so there is no solution.

Geometrically, the first two planes intersect in a line, and the plane $3x+y+3z=a$ is parallel to that line but does not contain it.
\end{enumerate}
\end{tcolorbox}


\newpage

% ---- Question 4 ----
\questionitem

The equations
\[
\begin{aligned}
x-y+2z&=4\\
(k+1)x+ky+(2-k)z&=4-k\\
(5-2k)x+(2-2k)y+(k+2)z&=8-k
\end{aligned}
\] \\
represent three planes, where $k$ is a real constant. \\
The planes do not intersect in exactly one point. \\
\begin{enumerate}[label=\textbf{(\alph*)}]
  \item Find the possible values of $k$. \marks{3} \\
  \item For each value of $k$ found in part (a), describe the configuration of the planes. \\
  Fully justify your answer, stating in each case whether the corresponding system is consistent or inconsistent. \marks{5} \\
\end{enumerate}

\vspace*{30pt}

\begin{tcolorbox}[colback=white, colframe=scarlet, title={\textbf{Solution}}, breakable, grow to left by=6mm, grow to right by=10mm]
\begin{enumerate}[label=\textbf{(\alph*)}]
  \item Let the coefficient matrix be
\[
A=\begin{pmatrix}
1 & -1 & 2\\
k+1 & k & 2-k\\
5-2k & 2-2k & k+2
\end{pmatrix}
\]

The planes fail to intersect in exactly one point when the system does \emph{not} have a unique solution, so we require
\[
\det A=0
\]

Expanding along the first row,
\begin{align*}
\det A
&=1\begin{vmatrix}k&2-k\\2-2k&k+2\end{vmatrix}
-(-1)\begin{vmatrix}k+1&2-k\\5-2k&k+2\end{vmatrix}
+2\begin{vmatrix}k+1&k\\5-2k&2-2k\end{vmatrix}\\
&=k(k+2)-(2-k)(2-2k)\\
&\quad +(k+1)(k+2)-(2-k)(5-2k)\\
&\quad +2\bigl((k+1)(2-2k)-k(5-2k)\bigr)\\
&=\bigl(k^2+2k-(4-6k+2k^2)\bigr)\\
&\quad +\bigl(k^2+3k+2-(10-9k+2k^2)\bigr)\\
&\quad +2\bigl((2-2k^2)-(5k-2k^2)\bigr)\\
&=(-k^2+8k-4)+(-k^2+12k-8)+(4-10k)\\
&=-2k^2+10k-8\\
&=-2(k^2-5k+4)\\
&=-2(k-1)(k-4)
\end{align*}

So
\[
-2(k-1)(k-4)=0
\]

Hence the possible values are
\[
k=1 \quad \text{or} \quad k=4
\]

\item We consider each value of \(k\) separately.

\medskip

For \(k=1\), the planes are
\[
\begin{aligned}
x-y+2z&=4\\
2x+y+z&=3\\
3x+3z&=7
\end{aligned}
\]

Adding the first two equations gives
\[
(x-y+2z)+(2x+y+z)=4+3
\]
so
\[
3x+3z=7,
\]
which is exactly the third equation.

Therefore any point that lies on the first two planes automatically lies on the third plane as well.

Now the normals of the first two planes are
\[
\mathbf n_1=(1,-1,2), \qquad \mathbf n_2=(2,1,1).
\]
These are not scalar multiples, so the first two planes are not parallel. Hence they intersect in a line, and the third plane contains this same line.

Therefore the system is consistent, and the three planes form a \textbf{sheaf}: they share a common line, so there are infinitely many solutions.

To find the line, add the first two equations:
\[
3x+3z=7 \quad \Rightarrow \quad x+z=\frac73.
\]
Then substitute into \(2x+y+z=3\):
\[
2x+y+\left(\frac73-x\right)=3
\]
so
\[
x+y=\frac23.
\]

Hence the common line is
\[
x+y=\frac23,\qquad x+z=\frac73.
\]

\medskip

For \(k=4\), the planes are
\[
\begin{aligned}
x-y+2z&=4\\
5x+4y-2z&=0\\
-3x-6y+6z&=4
\end{aligned}
\]

Suppose a point lies on the first two planes. Then
\[
2(x-y+2z)-(5x+4y-2z)=2\cdot 4-0,
\]
so
\[
-3x-6y+6z=8.
\]

But the third plane is
\[
-3x-6y+6z=4.
\]

This is impossible, so there is no point common to all three planes. Therefore the system is inconsistent.

Now check how the planes are arranged. Their normals are
\[
\mathbf n_1=(1,-1,2),\qquad
\mathbf n_2=(5,4,-2),\qquad
\mathbf n_3=(-3,-6,6).
\]
No two of these are scalar multiples, so no two planes are parallel. Hence each pair of planes meets in a line.

Let \(\mathbf d_{ij}\) be a direction vector for the line of intersection of planes \(i\) and \(j\). Then
\[
\mathbf d_{ij}=\mathbf n_i\times \mathbf n_j.
\]
Also,
\[
\mathbf n_3=2\mathbf n_1-\mathbf n_2.
\]

So
\[
\mathbf d_{13}
=\mathbf n_1\times \mathbf n_3
=\mathbf n_1\times(2\mathbf n_1-\mathbf n_2)
=-(\mathbf n_1\times \mathbf n_2)
=-\mathbf d_{12},
\]
and
\[
\mathbf d_{23}
=\mathbf n_2\times \mathbf n_3
=\mathbf n_2\times(2\mathbf n_1-\mathbf n_2)
=-2(\mathbf n_1\times \mathbf n_2)
=-2\mathbf d_{12}.
\]

Hence all three pairwise intersection lines are parallel.

They cannot all be the same line, because that would give a common line to all three planes, contradicting the fact that the system is inconsistent.

Therefore, when \(k=4\), the three planes form a \textbf{prism}: each pair of planes meets in a line, but the three lines are parallel and distinct, so there is no common point.
\end{enumerate} 
\end{tcolorbox} 

\newpage

% ---- Question 5 ----
\questionitem

Three planes have equations
\[
\begin{aligned}
x+ay&=-1\\
2x+y+z&=3\\
3x+by+2z&=c
\end{aligned}
\] \\
where $a$, $b$ and $c$ are constants. \\
\begin{enumerate}[label=\textbf{(\alph*)}]
  \item In the case where the planes \textbf{do not} intersect at a unique point,
  \begin{enumerate}[label=(\roman*)]
    \item find $b$ in terms of $a$ \marks{4}
    \item find the value of $c$ for which the planes form a sheaf. \marks{3} \\
  \end{enumerate}
  \item In the case where $b=1-a$ and $c=5$, find the coordinates of the point of intersection of the planes in terms of $a$. \marks{6} \\
\end{enumerate}

\vspace*{30pt}

\begin{tcolorbox}[colback=white, colframe=scarlet, title={\textbf{Solution}}, breakable, grow to left by=6mm, grow to right by=10mm]
\begin{enumerate}[label=\textbf{(\alph*)}]
  \item Let the planes be
  \[
  P_1:\ x+ay=-1,\qquad
  P_2:\ 2x+y+z=3,\qquad
  P_3:\ 3x+by+2z=c
  \]

  \begin{enumerate}[label=(\roman*)]
    \item For the three planes to intersect at a unique point, the coefficient matrix must have non-zero determinant:
    \[
    A=
    \begin{pmatrix}
    1 & a & 0\\
    2 & 1 & 1\\
    3 & b & 2
    \end{pmatrix}
    \]

    So
    \begin{align*}
    \det(A)
    &=
    \begin{vmatrix}
    1 & a & 0\\
    2 & 1 & 1\\
    3 & b & 2
    \end{vmatrix} \\
    &=1
    \begin{vmatrix}
    1 & 1\\
    b & 2
    \end{vmatrix}
    -a
    \begin{vmatrix}
    2 & 1\\
    3 & 2
    \end{vmatrix}
    +0
    \begin{vmatrix}
    2 & 1\\
    3 & b
    \end{vmatrix} \\
    &=1(2-b)-a(4-3) \\
    &=2-b-a
    \end{align*}

    In the case where the planes do \textbf{not} intersect at a unique point,
    \[
    \det(A)=0
    \]
    so
    \begin{align*}
    2-b-a&=0\\
    b&=2-a
    \end{align*}

    Hence \(b=2-a\).

    \item A sheaf of planes means that all three planes contain a common line.

    First find the line of intersection of \(P_1\) and \(P_2\).

    From \(P_1\),
    \[
    x=-1-ay
    \]

    Substitute into \(P_2\):
    \begin{align*}
    2x+y+z&=3\\
    2(-1-ay)+y+z&=3\\
    -2-2ay+y+z&=3\\
    z&=5+(2a-1)y
    \end{align*}

    So the line \(P_1\cap P_2\) has coordinates
    \[
    x=-1-ay,\qquad z=5+(2a-1)y
    \]

    For \(P_3\) to contain this whole line, substitute these expressions into \(P_3\), using \(b=2-a\):
    \begin{align*}
    3x+(2-a)y+2z
    &=3(-1-ay)+(2-a)y+2\bigl(5+(2a-1)y\bigr)\\
    &=-3-3ay+(2-a)y+10+(4a-2)y\\
    &=7+\bigl(-3a+2-a+4a-2\bigr)y\\
    &=7
    \end{align*}

    This is satisfied for every point on the line precisely when
    \[
    c=7
    \]

    Hence \(c=7\).
  \end{enumerate}

  \item Now take \(b=1-a\) and \(c=5\).

  From the first plane,
  \[
  x=-1-ay
  \]

  Substitute into the second plane:
  \begin{align*}
  2x+y+z&=3\\
  2(-1-ay)+y+z&=3\\
  -2-2ay+y+z&=3\\
  z&=5+(2a-1)y
  \end{align*}

  Now substitute \(x=-1-ay\) and \(z=5+(2a-1)y\) into the third plane
  \[
  3x+(1-a)y+2z=5
  \]

  Then
  \begin{align*}
  3(-1-ay)+(1-a)y+2\bigl(5+(2a-1)y\bigr)&=5\\
  -3-3ay+(1-a)y+10+(4a-2)y&=5\\
  7+\bigl(-3a+1-a+4a-2\bigr)y&=5\\
  7-y&=5\\
  y&=2
  \end{align*}

  Now find \(x\) and \(z\):
  \begin{align*}
  x&=-1-a(2)=-1-2a\\
  z&=5+(2a-1)(2)=5+4a-2=4a+3
  \end{align*}

  Therefore the point of intersection is
  \[
  (-1-2a,\ 2,\ 4a+3)
  \]

  Hence the coordinates are \((-1-2a,\,2,\,4a+3)\).
\end{enumerate}
\end{tcolorbox}


\newpage

% ---- Question 6 ----
\questionitem

\[
\mathbf{M}=\begin{pmatrix}
1&1&0\\
3&k&1\\
2&-1&1
\end{pmatrix}
\] \\
where $k$ is a constant. \\
\begin{enumerate}[label=\textbf{(\alph*)}]
  \item Find the values of $k$ for which the matrix $\mathbf{M}$ has an inverse. \marks{2} \\
  \item Find, in terms of $p$, the coordinates of the point where the following planes intersect:
  \[
  \begin{aligned}
  x+y&=1\\
  3x+2y+z&=p\\
  2x-y+z&=4
  \end{aligned}
  \]
  \marks[-20pt]{5}
  \item
  \begin{enumerate}[label=(\roman*)]
    \item Find the value of $q$ for which the set of simultaneous equations
    \[
    \begin{aligned}
    x+y&=1\\
    3x+z&=q\\
    2x-y+z&=4
    \end{aligned}
    \]
    has at least one solution. \marks[-20pt]{2}
    \item For this value of $q$, interpret the solution set geometrically. \marks{2} \\
  \end{enumerate}
\end{enumerate}

\vspace*{30pt}

\begin{tcolorbox}[colback=white, colframe=scarlet, title={\textbf{Solution}}, breakable, grow to left by=6mm, grow to right by=10mm]
\begin{enumerate}[label=\textbf{(\alph*)}]
  \item A matrix has an inverse if and only if its determinant is non-zero.

  For
  \[
  \mathbf{M}=\begin{pmatrix}
  1 & 1 & 0\\
  3 & k & 1\\
  2 & -1 & 1
  \end{pmatrix}
  \]
  expand along the first row:
  \begin{align*}
  \det(\mathbf{M})
  &=1\begin{vmatrix}k&1\\-1&1\end{vmatrix}
   -1\begin{vmatrix}3&1\\2&1\end{vmatrix}
   +0\begin{vmatrix}3&k\\2&-1\end{vmatrix} \\
  &=1(k\cdot 1-1\cdot(-1)) -1(3\cdot 1-1\cdot 2) +0 \\
  &=(k+1)-(3-2) \\
  &=k
  \end{align*}

  Therefore $\mathbf{M}$ has an inverse when
  \[
  k\neq 0
  \]

  \item We solve the three plane equations
  \[
  \begin{aligned}
  x+y&=1 \\
  3x+2y+z&=p \\
  2x-y+z&=4
  \end{aligned}
  \]

  From the first equation,
  \[
  y=1-x
  \]

  Substitute into the second equation:
  \begin{align*}
  3x+2(1-x)+z&=p \\
  3x+2-2x+z&=p \\
  x+z&=p-2
  \end{align*}

  Substitute into the third equation:
  \begin{align*}
  2x-(1-x)+z&=4 \\
  2x-1+x+z&=4 \\
  3x+z&=5
  \end{align*}

  So we now have
  \[
  x+z=p-2
  \qquad\text{and}\qquad
  3x+z=5
  \]

  Subtract the first from the second:
  \begin{align*}
  (3x+z)-(x+z)&=5-(p-2) \\
  2x&=7-p \\
  x&=\frac{7-p}{2}
  \end{align*}

  Then
  \begin{align*}
  y&=1-x \\
   &=1-\frac{7-p}{2} \\
   &=\frac{2-7+p}{2} \\
   &=\frac{p-5}{2}
  \end{align*}

  And using $3x+z=5$,
  \begin{align*}
  z&=5-3x \\
   &=5-3\left(\frac{7-p}{2}\right) \\
   &=\frac{10-21+3p}{2} \\
   &=\frac{3p-11}{2}
  \end{align*}

  Hence the planes intersect at
  \[
  \left(\frac{7-p}{2},\,\frac{p-5}{2},\,\frac{3p-11}{2}\right)
  \]

  \item
  \begin{enumerate}[label=(\roman*)]
    \item The equations are
    \[
    \begin{aligned}
    x+y&=1 \\
    3x+z&=q \\
    2x-y+z&=4
    \end{aligned}
    \]

    Using the first two equations,
    \[
    (3x+z)-(x+y)=q-1
    \]

    But
    \[
    (3x+z)-(x+y)=2x-y+z
    \]

    So the third equation requires
    \[
    q-1=4
    \]

    Hence
    \[
    q=5
    \]

    \item For $q=5$, the second equation is
    \[
    3x+z=5
    \]

    Then
    \[
    (3x+z)-(x+y)=5-1=4
    \]
    which gives
    \[
    2x-y+z=4
    \]

    So the third equation is implied by the first two, and does not give any new restriction.

    Therefore the three planes have infinitely many common points, forming a line.

    Let $x=t$. Then
    \[
    y=1-t
    \]
    and from $3x+z=5$,
    \[
    z=5-3t
    \]

    So the solution set is
    \[
    (x,y,z)=\bigl(t,\,1-t,\,5-3t\bigr),\qquad t\in\mathbb{R}
    \]

    Geometrically, the three planes intersect in a common line, so form a \textbf{sheaf}.
  \end{enumerate}
\end{enumerate}
\end{tcolorbox}


\newpage

% ---- Question 7 ----
\questionitem

You are given that $a$ is a parameter which can take only real values. \\
The matrix $A$ is given by
\[
A=\begin{pmatrix}
3&2&a\\
1&3&0\\
2&-1&1
\end{pmatrix}
\] \\
\begin{enumerate}[label=\textbf{(\alph*)}]
  \item Find an expression for the determinant of $A$ in terms of $a$. \marks{2} \\
\end{enumerate}
You are given the following system of equations in $x$, $y$ and $z$.
\[
\begin{aligned}
3x+2y+az&=8\\
x+3y&=5\\
2x-y+z&=3
\end{aligned}
\] \\
The system can be written in the form
\[
A\begin{pmatrix}x\\y\\z\end{pmatrix}=\begin{pmatrix}8\\5\\3\end{pmatrix}
\] \\
\begin{enumerate}[label=\textbf{(\alph*)}, resume]
  \item
  \begin{enumerate}[label=(\roman*)]
    \item In the case where $A$ is not singular, solve the given system by using $A^{-1}$. \marks{5}
    \item In the case where $A$ is singular, describe the configuration of the planes whose equations are the three equations of the system. \marks{3} \\
  \end{enumerate}
\end{enumerate}
The transformation represented by $A$ is denoted by $T$. \\
A 3-D solid of volume $a^2+3$ is transformed by $T$ to an image. \\
\begin{enumerate}[label=\textbf{(\alph*)}, resume]
  \item
  \begin{enumerate}[label=(\roman*)]
    \item Determine the range of values of $a$ for which the orientation of the image is the reverse of the orientation of the object. \marks{1}
    \item Determine the range of values of $a$ for which the volume of the image is less than the volume of the object. \marks{2} \\
  \end{enumerate}
\end{enumerate}

\vspace*{30pt}

\begin{tcolorbox}[colback=white, colframe=scarlet, title={\textbf{Solution}}, breakable, grow to left by=6mm, grow to right by=10mm]
\begin{enumerate}[label=\textbf{(\alph*)}]
  \item Expanding $\det(A)$ along the first row,
\begin{align*}
\det(A)
&= 3\begin{vmatrix}3&0\\-1&1\end{vmatrix}
-2\begin{vmatrix}1&0\\2&1\end{vmatrix}
+a\begin{vmatrix}1&3\\2&-1\end{vmatrix} \\
&= 3(3\cdot 1-0\cdot(-1))
-2(1\cdot 1-0\cdot 2)
+a(1\cdot(-1)-3\cdot 2) \\
&= 3(3)-2(1)+a(-1-6) \\
&= 9-2-7a \\
&= 7(1-a)
\end{align*}

Hence,
\[
\det(A)=7(1-a)
\]

  \item \begin{enumerate}[label=(\roman*)]
    \item For $A$ to be non-singular, we need
\[
\det(A)\neq 0
\]
So
\[
7(1-a)\neq 0 \implies a\neq 1
\]

For $a\neq 1$,
\[
A^{-1}=\frac{1}{7(1-a)}
\begin{pmatrix}
3 & -a-2 & -3a \\
-1 & 3-2a & a \\
-7 & 7 & 7
\end{pmatrix}
\]

Now
\[
\begin{pmatrix}x\\y\\z\end{pmatrix}
=A^{-1}\begin{pmatrix}8\\5\\3\end{pmatrix}
\]

So
\begin{align*}
\begin{pmatrix}x\\y\\z\end{pmatrix}
&=
\frac{1}{7(1-a)}
\begin{pmatrix}
3 & -a-2 & -3a \\
-1 & 3-2a & a \\
-7 & 7 & 7
\end{pmatrix}
\begin{pmatrix}8\\5\\3\end{pmatrix} \\
&=
\frac{1}{7(1-a)}
\begin{pmatrix}
3(8)+(-a-2)(5)+(-3a)(3) \\
(-1)(8)+(3-2a)(5)+a(3) \\
(-7)(8)+7(5)+7(3)
\end{pmatrix} \\
&=
\frac{1}{7(1-a)}
\begin{pmatrix}
24-5a-10-9a \\
-8+15-10a+3a \\
-56+35+21
\end{pmatrix} \\
&=
\frac{1}{7(1-a)}
\begin{pmatrix}
14-14a \\
7-7a \\
0
\end{pmatrix} \\
&=
\begin{pmatrix}
2\\1\\0
\end{pmatrix}
\end{align*}

Hence, for $a\neq 1$,
\[
x=2,\quad y=1,\quad z=0
\]

    \item $A$ is singular when
\[
7(1-a)=0 \implies a=1
\]

In this case the equations are
\begin{align*}
3x+2y+z&=8 \\
x+3y&=5 \\
2x-y+z&=3
\end{align*}

Now add the second and third equations:
\begin{align*}
(x+3y)+(2x-y+z)&=5+3 \\
3x+2y+z&=8
\end{align*}

This is exactly the first equation.

So the first plane contains the line of intersection of the second and third planes.

Also, the second and third planes are not parallel, since their normal vectors
\[
(1,3,0)\quad \text{and} \quad (2,-1,1)
\]
are not scalar multiples of each other.

Therefore the second and third planes intersect in a line, and the first plane also contains this same line.

Hence the configuration is: the three distinct planes all share a common line, so form a \textbf{sheaf}.
  \end{enumerate}

  \item \begin{enumerate}[label=(\roman*)]
    \item Orientation is reversed when
\[
\det(A)<0
\]

Since
\[
\det(A)=7(1-a)
\]
we need
\[
7(1-a)<0 \implies 1-a<0 \implies a>1
\]

Hence, the orientation is reversed for
\[
a>1
\]

    \item Volume scale factor under the transformation $T$ is
\[
|\det(A)|=|7(1-a)|
\]

So the image volume is
\[
|7(1-a)|(a^2+3)
\]

We want this to be less than the original volume $a^2+3$:
\[
|7(1-a)|(a^2+3)<a^2+3
\]

Since
\[
a^2+3>0
\]
for all real $a$, divide through by $a^2+3$ to get
\[
|7(1-a)|<1
\]

So
\[
7|1-a|<1 \implies |1-a|<\frac{1}{7}
\]

Hence
\[
-\frac{1}{7}<1-a<\frac{1}{7}
\]

Subtracting $1$ throughout gives
\[
-\frac{8}{7}<-a<-\frac{6}{7}
\]

Multiplying through by $-1$ and reversing the inequalities,
\[
\frac{6}{7}<a<\frac{8}{7}
\]

Hence, the image volume is less than the object volume for
\[
\frac{6}{7}<a<\frac{8}{7}
\]
  \end{enumerate}
\end{enumerate}
\end{tcolorbox}


\end{enumerate}
\end{document}
